\documentclass{ltjsarticle}
\usepackage{tcolorbox,mathcomp,tcolorbox,amsmath,mathtools,amssymb,graphicx,ascmac,fancybox,framed,tikz,comment,calc,enumitem}
\begin{document}
%
　定数$b,c,p,q,r$に対し，$x^4+bx+c=(x^2+px+q)(x^2-px+r)$が$x$に
の恒等式であるとする。
\begin{enumerate}
\item[(1)]\quad$p\ne{0}$であるとき，$q,r$を$p,b$で表せ。
\item[(2)]\quad$p\ne{0}$とする。$b,c$が定数$a$を用いて$b=(a^2+1)(a+2)$，$c=-\left(a+\dfrac{3}{4}\right)(a^2+1)$
と表されているとき，有理数を係数とする$t$についての整式$f(t)$と$g(t)$で
\[\{p^2-(a^2+1)\}\{p^4+f(a)p+g(a)\}=0\]
を満たすものを$1$組求めよ。
\item[(3)]\quad$a$を整数とする。$x$の$4$次式$x^4+(a^2+1)(a+2)x-\left(a+\dfrac{3}{4}\right)(a^2+1)$が有理数を係数とする$2$次式に因数分解できるような$a$をすべて求めよ。
\end{enumerate}
%
\end{document}
